% Chapter 5: Testing and Quality with AI-Generated Code
\chapter{Testing and Quality with AI-Generated Code}
\label{ch:testing}

AI-generated code has a specific failure mode that human-written code does not:
it looks correct. The syntax is clean, the variable names are reasonable,
the logic reads well. And yet it may be subtly wrong in ways that only
manifest at runtime, under load, or at edge cases.

Testing is the antidote.

\section{Verification Is King}
\label{sec:verification}

\convergence{``Give Claude a way to verify its work'' is Anthropic's stated
single highest-leverage practice.\sourceurl{https://code.claude.com/docs/en/best-practices}
This aligns exactly with what practitioners discover independently:
the projects that succeed are the ones with verification at every layer.}

\subsection{The 6-Layer Validation Architecture}

Layer defenses so that no single failure mode goes undetected:

\begin{enumerate}
  \item \textbf{Type Safety} --- Catch type errors at compile/lint time,
    not runtime. TypeScript strict mode, Python type hints with mypy,
    Rust's borrow checker.

  \item \textbf{Input Validation} --- Check preconditions at function entry.
    Fail fast with explicit errors. Never let invalid data propagate silently.

  \item \textbf{Unit Tests} --- Test each function in isolation.
    Happy path + error cases + edge cases.
    Target: 80\%+ coverage for modules.

  \item \textbf{Integration Tests} --- Test multi-function workflows.
    Verify components work together with realistic data.

  \item \textbf{End-to-End Tests} --- Test complete user workflows
    from input to output. Verify the system meets requirements.

  \item \textbf{Property-Based Tests} --- Test invariants that should always hold.
    Generate random inputs. Catch edge cases you did not think of.
\end{enumerate}

\margintip{Layers 1--2 are cheap to add to any project today. Layers 3--4 are the baseline for production. Layers 5--6 are for critical systems.}

\subsection{The Pattern to Avoid}

\begin{warnbox}[Cautionary Tale]
A common failure mode: a developer generates four modules totaling 1,500+
lines of code with zero tests. The code looks clean. It compiles. It appears to work.

Result: zero confidence in functionality, impossible to refactor,
and months of accumulated technical debt.

This pattern is common enough to have a name: the Verification Gap
(see \cref{ch:antipatterns}).
\end{warnbox}

% -------------------------------------------------------------------
\section{Phase-Appropriate Standards}
\label{sec:phases}

\practitioner{Not all code needs the same rigor.
Applying production standards to a prototype kills velocity.
Applying prototype standards to production kills reliability.}

\subsection{Three Development Phases}

\begin{fullwidth}
\begin{tabular}{@{}lp{3cm}p{3cm}p{4cm}@{}}
\toprule
\textbf{Phase} & \textbf{Testing} & \textbf{Code Quality} & \textbf{Transition Criteria} \\
\midrule
Exploration & Manual verification OK & Long functions OK, mutable state OK & Hypothesis validated; proceeding to build \\
Development & Unit + integration tests required & Style enforcement, type hints & Tests pass, coverage >60\%, code review \\
Production & Full 6-layer validation & Strict linting, immutability default & Coverage >80\%, zero critical warnings \\
\bottomrule
\end{tabular}
\end{fullwidth}

\subsection{Phase Transitions}

\practitioner{Make phase transitions explicit.} Write them in CLAUDE.md:

\begin{minted}{markdown}
## Development Phase: EXPLORATION
- Tests: manual OK for now
- Coverage target: none
- Graduation criteria:
  - [ ] Core algorithm validated on real data
  - [ ] API surface finalized
  - [ ] Dependencies locked
  When all checked: switch to DEVELOPMENT phase
\end{minted}

The key insight: premature rigor kills exploration,
but sloppy production kills credibility.
The solution is not to choose one standard---it is to be explicit about
which standard applies \emph{right now} and when to transition.

\subsection{Verification Criteria in CLAUDE.md}

\convergence{Set verification expectations from day one.}
Even in the exploration phase, CLAUDE.md should contain:

\begin{minted}{markdown}
## Verification
- Always run tests after code changes
- Never commit without passing CI
- Include test plan in PR description
\end{minted}

These three lines improve every interaction by giving Claude
a concrete way to validate its own work.

\begin{keyinsight}
Testing AI-generated code is not about distrust---it is about
leveraging the AI's own capabilities. Claude writes excellent tests
when given clear instructions. The developer's job is to define
\emph{what} to test; Claude handles the \emph{how}.
\end{keyinsight}
